% Phase 2 Hybrid ML + DL paragraph for IEEE report
% Author: Daksh Batra

\subsection{Hybrid ML and Deep Learning Feature Fusion}
While handcrafted photometric and texture descriptors achieved strong baseline classification in Phase~1 ($0.8803\pm0.0215$ mean cross-validation macro F1 across five folds), systematic confusions persisted among glass, metal, and plastic due to overlapping specular reflections and smooth surface gradients. To investigate whether high-level geometric representations resolve these material confusions, we evaluated a hybrid architecture combining off-the-shelf deep convolutional neural network (CNN) representations~\cite{razavian2014cnn} with engineered descriptors. Pretrained 1,280-dimensional feature embeddings from frozen MobileNetV2 and EfficientNet-B0 backbones were extracted and fused with Phase~1's winning 1,608-dimensional extended color and spatial texture features. To prevent the higher-dimensional handcrafted block from dominating the radial basis function (RBF) kernel distance metric, we enforced a fold-local dimensional normalization protocol: each block was independently standardized using training-fold statistics and scaled by $1/\sqrt{D_{\mathrm{block}}}$ prior to concatenation ($D_{\mathrm{total}} = 2,888$), giving both modalities equal total geometric variance. Evaluating across the five frozen training cross-validation folds, early feature fusion with an RBF SVM ($C=10.0$, $\gamma=\text{`scale'}$) achieved a mean cross-validation macro F1 of $0.9044\pm0.0332$ and an overall accuracy of 91.73\% (compared to $0.8597\pm0.0194$ for classical classifiers trained on CNN embeddings alone). 

In accordance with our pre-registered decision rule ($\Delta > \sigma$), the $+0.0241$ absolute gain of early fusion under initial discrete predictions ($0.9044\pm0.0332$ Macro F1) did not constitute a clear win under the pre-registered rule as it remained within the cross-fold sample standard deviation ($\sigma = 0.0332$); when evaluated under the argmax of predicted probabilities for strict cross-model consistency, H2 achieves $0.9175\pm0.0262$ with 92.52\% pooled accuracy (1,632/1,764). Following the selection of the fine-tuned MobileNetV2 architecture (A2.1, $0.8825\pm0.0219$ macro F1, 89.80\% pooled accuracy) in Phase~2 ablations, we evaluated late probability fusion (H3) combining out-of-fold probability vectors from A2.1 and the calibrated Phase~1 SVM via convex blending ($p = w \cdot p_{\mathrm{cnn}} + (1 - w) \cdot p_{\mathrm{svm}}$). Analysis of out-of-fold predictions reveals that while CNN and SVM predictions disagree on 16.72\% of training images (295/1,764), both models fail simultaneously on only 2.61\% (46/1,764), demonstrating complementary error patterns. Using nested cross-validation across the five frozen training folds—where fusion weight $w \in [0, 1]$ was selected on inner training folds independently for each outer fold—equal weighting ($w=0.5$) was selected across all folds. H3 achieved a mean cross-validation macro F1 of $\mathbf{0.9364\pm0.0171}$, mean per-fold accuracy of 94.62\%, and pooled accuracy of $\mathbf{94.61\%}$ (1,669/1,764). Comparing against the best standalone single model (A2.1, 0.8825 macro F1), this represents a clear win under the pre-registered rule ($\Delta = 0.0539 > \sigma = 0.0219$). When compared to H2, $\Delta = 0.0189 < \sigma_{\mathrm{H2}} = 0.0262$, so H3 is carried forward as the highest mean candidate per protocol rather than a clearly separated winner under the rule over H2. All cross-validation estimates remain somewhat optimistic as base models were selected on these five folds; the final test split ($N=377$) provides the unbiased assessment. Pooled confusion analysis across all 1,764 training samples demonstrated that metal recall reached 98.25\% (281/286), reducing metal-to-glass errors to just 3 (an 85.7\% reduction compared to Phase~1), while plastic recall reached 91.67\% (308/336) and trash recall rose to 86.46\% (83/96). The final hybrid model was assembled on the full training split—preserving the verified A2.1 checkpoint (SHA-256: \texttt{26592f8b...5d28e}) without retraining—and exported for downstream probability calibration and operational abstention threshold selection on the held-out validation split.
